% 日本語版（研究者が読むためのもの）．英語版 ../main.tex と節・図・表・式・数値を 1 対 1 に対応させる．
% ビルドは build.sh（LuaLaTeX と BibTeX）．数値は英語版と同じく \NUM{key} で書き，../numbers.tex を読む．
\documentclass[paper=a4,fontsize=9.5pt,twocolumn,line_length=26zw,number_of_lines=50,head_space=17mm,baselineskip=1.55zw,column_gap=2zw]{jlreq}
\usepackage[ipaex]{luatexja-preset}
\usepackage{amsmath,amssymb}
\usepackage{amsthm}
\usepackage{graphicx}
\usepackage{booktabs}
\usepackage{multirow}
\usepackage[nocompress]{cite}
\usepackage{url}
\usepackage[hidelinks]{hyperref}

% 命題の環境（英語版の \newtheorem{proposition}{Proposition} に対応）．本文は立体にする（IEEEtran と同じ）．
\newtheoremstyle{japroposition}{0.5\baselineskip}{0.5\baselineskip}{\normalfont}{0pt}{\sffamily\gtfamily\bfseries}{}{1\zw}{\thmname{#1}\thmnumber{#2}\thmnote{（#3）}}
\theoremstyle{japroposition}
\newtheorem{proposition}{命題}

% \emph の中の和文はゴシックにする（欧文は斜体のまま）．
\let\jaorigemph\emph
\renewcommand{\emph}[1]{\jaorigemph{\gtfamily #1}}

% 本文の実験由来の数値は \NUM{key} で書く（英語版と同じ台帳 ../results/numbers.json → ../numbers.tex）．
% 自動生成．手で編集しない．
% 生成元: results/numbers.json ／ 生成器: scripts/gen_numbers_tex.py
% 本文の実験由来の数値は \NUM{key} で書く．未登録の key はビルド時にエラーになる．
\makeatletter
\providecommand{\NUM}[1]{%
  \@ifundefined{NUM@#1}{%
    \GenericError{}{numbers.tex: key '#1' is not registered in results/numbers.json}{}{}%
    \textbf{??}%
  }{\csname NUM@#1\endcsname}%
}
\makeatother


% 図：英語版の \FIG と同じ．図は ../figures/<name>.pdf を使う．
\newcommand{\FIG}[2]{\IfFileExists{../figures/#1.pdf}{\includegraphics[width=#2]{../figures/#1.pdf}}{\fbox{\parbox[c][6em][c]{\dimexpr#2-2\fboxsep-2\fboxrule\relax}{\centering\footnotesize figure \texttt{#1} (pending)}}}}

\renewcommand{\abstractname}{要旨}

\title{少数のターゲットラベルが最初に直すもの：\\コーパス横断のマルチモーダル感情認識における\\判断基準と向きの分析，\\およびモダリティをまたぐ弱いラベルの転移}

\author{著者名（未記入）\thanks{所属（未記入）．}}

\date{IEEE Transactions on Affective Computing 投稿原稿の日本語版}

% jlreq では，2 段組のときにプリアンブルに書いた abstract が \maketitle で全幅に置かれる．
\begin{abstract}
新しいコーパスに移した認識器は，感情ごとに少数のラベルと，他のコーパスのラベル付きデータを持つことが多い．
本論文は，凍結した（微調整しない）音声・顔・テキストの特徴の上の線形分類器において，これらのラベルが何を直すかを，7つのコーパス群のすべての順序対にわたって分析する．
感情の対ごとの判定を判断基準（threshold）と向き（direction）に分けると，単一モダリティの音声と顔で感情ごとのラベルが5本以下のとき，ラベルが合併（pooled）の分類器に加える改善のうち\NUM{p1-e29-critshare-kle5-single-lo}から\NUM{p1-e29-critshare-kle5-single-hi}の割合が判断基準の修正から生じる．
一方，音声の向きは\NUM{p1-qdir-audio-lo}から\NUM{p1-qdir-audio-hi}の実効次元（effective dimension）で決まる速さで改善し，1つのソースコーパスは向きについて感情ごとに\NUM{p1-e32-n1src-audio-lo}から\NUM{p1-e32-n1src-audio-hi}本のラベルに相当する．
少数のラベルも，ラベルを使わない選択も，強い出発点（start）の向きを直さない．
そこで本論文は，ソースだけで学習し，生徒（student）がもたないモダリティを見る融合の分類器が，ラベルのないターゲットの発話すべてに付ける弱いラベル（weak labels）で，出発点にアンカーした線形の生徒を学習する．
ターゲットのプールで学び直した先生を加えて拡張すると，この手法は，感情ごとのラベルが0本と5本のとき，第2のエンコーダの生徒を含む11の生徒すべてについて，同じ方法で調整した9つの比較手法のうち最も強いものをUARで\NUM{cmp-oursx-vs-best-range-lo}から\NUM{cmp-oursx-vs-best-range-hi}ポイント上回る．

\medskip
\noindent{\sffamily\gtfamily\bfseries 索引語}\quad
マルチモーダル感情認識，コーパス横断，少数ショット学習，ドメイン適応，知識蒸留，弱いラベル．
\end{abstract}

% 段の下端をそろえるために式の前後が引き伸ばされるのを防ぐ．
\raggedbottom
% 図表だけの頁では，図表を頁の上から詰めて置く．
\makeatletter
\setlength{\@fptop}{0pt}\setlength{\@fpsep}{\floatsep}\setlength{\@fpbot}{0pt plus 1fil}
\setlength{\@dblfptop}{0pt}\setlength{\@dblfpsep}{\dblfloatsep}\setlength{\@dblfpbot}{0pt plus 1fil}
\makeatother
% 図と表の説明文：ラベル（図1，表1）はゴシックの太字，本文は明朝で small．
\jlreqsetup{caption_font={\small},caption_label_font={\small\sffamily\gtfamily\bfseries}}

\begin{document}

\maketitle

\section{はじめに}
\label{sec:introduction}

新しいコーパスに配備した感情認識器は，学習データに含まれなかった話者，録音条件，感情の引き出し方，注釈の慣習に出会い，コーパスをまたぐと正解率が下がる\cite{schuller2010crosscorpus,gideon2021addog,feng2023fewshot}．
新しいコーパスにラベルを付けるには費用がかかるので，実務家は通常，感情ごとに少数の発話にラベルを付け，他のコーパスのラベル付きの発話を加える．
それらのソースコーパスが役立つか，どれだけ役立つかは，コーパスとターゲットのラベルの本数によって決まる．
マルチモーダル感情認識では，ターゲットの発話が複数のモダリティを伴うことが多く，コーパスをまたいで学習する学習器はこれを第2の情報源として使える．

この設定を扱う研究には2つの系統があり，それぞれ別の面から取り組んでいる．
複数のドメインからの学習の理論は，ソースを合併するとターゲットのラベルが少ないときは役立ち，多いときは害になることを示し，ターゲットの誤差を，ソースの誤差，分布の隔たり（divergence），ターゲットの標本の大きさで上から抑える\cite{bendavid2010theory,blitzer2007learning,desilva2023value}．
少数のターゲットラベルによる感情認識の手法は，音声のモデルや音声とテキストのモデルを少数のラベル付き発話で適応させ\cite{feng2023fewshot,he2025m4ser,zaidi2024enhancing}，半教師あり学習と蒸留は，ラベルのないデータやより強い先生（teacher）を使う\cite{lee2013pseudolabel,hinton2015distilling,lopezpaz2016unifying}．
少数のターゲットラベルによるコーパス横断の設定において，どちらの系統も，少数のターゲットラベルが分類器のどの部分を直すか，ソースがラベルの代わりにどれだけを補うか，配備したモデルが1つのモダリティしか見ないときに第2のモダリティが何を加えうるかを示していない．
顔を含む複数のモダリティを使うコーパス横断の手法も，依然として少ない\cite{ahn2025multitask}．

本論文は，凍結した（微調整しない）エンコーダの特徴の上の線形ヘッドについて，3つの問いに答える．
この設定は分析できるほど単純であり，4つのターゲットで音声エンコーダを低ランクの適応で微調整した確認では，微調整より強い．
(Q1)~少数のターゲットラベルは，分類器の何を直すか．
(Q2)~ソースは，ラベルの代わりにどれだけを補うか．
(Q3)~その答えから，同じ手順で調整した標準的な半教師あり学習と蒸留の手法を上回る手法が得られるか．
6つのコーパスから作った7つの群を，音声・顔・テキストの特徴とそれらの融合で調べる（図\ref{fig:overview}）．
4クラス分類器の感情の対ごとの判定を判断基準のずれ（threshold offset）と識別力（separation）で表し，それぞれについて感情ごとに$K$本のラベルでのコストを導き，その予言を，コーパスのすべての順序対の利得の地図（gain map）と照らし合わせる．

\begin{figure*}[t]
\centering
\FIG{overview}{\textwidth}
\caption{概観．出発点の分類器は，凍結したエンコーダを用いて，ラベル付きのソースコーパスで学習する．第I部：単一モダリティの音声と顔の特徴で$K \le 5$のとき，感情ごとに$K$本のラベルが加える改善の大半は判断基準から生じ（残る誤差は$1/(2K)$），音声の向きはゆっくり改善する（残る誤差は$q/K$）．第II部：ソースだけで学習し，生徒がもたないモダリティを含む融合の分類器が，ラベルのないターゲットの発話すべてに，感情ごとの合計の重みが等しい弱いラベルを付け，出発点にアンカーした線形の生徒が，この弱いラベルと$K$本のラベルから学ぶ．テストのとき，生徒は自分のモダリティだけを必要とする．}
\label{fig:overview}
\end{figure*}

本論文の貢献は次のとおりである．
\begin{enumerate}
\item 顔を含む3つのモダリティとその融合について，7つの群のすべての順序対と，感情ごとに1本からプール全体までのラベルの本数にわたって作った，コーパス横断の合併の利得の地図．30回の反復にもとづく対応のある区間を含む（\ref{sec:part1-map}節）．
\item 利得の判断基準の部分と向きの部分への分解．分解は$K$の関数として測り，同じ分割の上でラベルを使わない中心化と比べる．単一モダリティの音声と顔で$K \le 5$のとき，ラベルが合併の分類器に加える改善の\NUM{p1-e29-critshare-kle5-single-lo}--\NUM{p1-e29-critshare-kle5-single-hi}が判断基準から生じ，音声の分類器の向きは$q/K$の速さでしか改善しない（測った$q$は\NUM{p1-qdir-audio-lo}--\NUM{p1-qdir-audio-hi}）（\ref{sec:part1-fix}節）．
\item ソースの価値をターゲットのラベルの本数に換算する式$n = q e^2/(1-e^2)$．この式は推定した向きと標本サイズの古典的な関係を逆に解いて得たものであり，測った値は，音声のソース1つで感情ごとに\NUM{p1-e32-n1src-audio-lo}--\NUM{p1-e32-n1src-audio-hi}本，音声のソースすべてを合併すると\NUM{p1-e32-nall-audio-lo}--\NUM{p1-e32-nall-audio-hi}本である（\ref{sec:part1-value}節）．
\item ソースだけで学習し，ターゲットのデータで学び直さず，生徒がもたないモダリティを見る融合の分類器から，クラスの釣り合いをとった弱いラベルをラベルのないターゲットのプール全体に付け，仮のターゲットで$K$ごとに選んだハイパーパラメータで，アンカーつきの線形の生徒を学習する手法．基本の手法は，生徒と$K \in \{0, 5\}$の\NUM{cmp-ours-vs-best-cells-n}の組み合わせのうち\NUM{cmp-ours-vs-best-pos-n}で，同じ手順で調整した最も強い比較手法を上回る．ターゲットのプールで学び直した先生を加え，弱いラベルの重みづけを$K$ごとに選ぶ拡張版は，そのすべての組み合わせで最も強い比較手法を\NUM{cmp-oursx-vs-best-range-lo}--\NUM{cmp-oursx-vs-best-range-hi}ポイント上回り，第2のエンコーダでも上回る．本論文で測った場合には，先生の出発点に対するUARの優位が利得を上から抑える（\ref{sec:part2-crossmodal}節と\ref{sec:comparison}節）．
\end{enumerate}
我々の知る限り，これら4つの設定と測定の組み合わせはこれまで報告されていない．
\ref{sec:related}節では，各構成要素と各式を先行研究に対して位置づける．

\ref{sec:related}節で関連研究を概観し，\ref{sec:setting}節で手順を定義し，\ref{sec:part1-what-few-labels-fix}節で第I部，\ref{sec:part2-mechanism-to-method}節で第II部，\ref{sec:comparison}節で比較を示す．
\ref{sec:discussion}節と\ref{sec:limitations-ethics}節では，結果の意味と限界を論じる．

\section{関連研究}
\label{sec:related}

\emph{少数のターゲットラベルによるコーパス横断の感情認識．}
コーパス横断の評価は，音声感情認識で大きな性能の低下を以前から示しており，コーパスごとの正規化の動機となってきた\cite{schuller2010crosscorpus}．
少数のラベル付きターゲット発話を使う手法には，自発的な音声のためのシャムネットワークによる少数ショット学習\cite{feng2023fewshot}，ターゲットのラベルの量を増やしながら評価した敵対的なドメイン汎化\cite{gideon2021addog}，ソースコーパスでの事前の微調整（pre-finetuning）\cite{chen2023prefinetuning}，演技された音声・映像と即興の音声・映像の間の適応\cite{li2021acted}，凍結した音声特徴の上でロジスティック回帰を用いる能動学習\cite{xu2024crosstask}がある．
音声とテキストのモデルは，クラスごとに5から60本のラベルで\cite{he2025m4ser}，あるいは言語をまたいでコーパスごとに少数のラベルで\cite{zaidi2024enhancing}適応されている．
これらの研究は手法を提案するか評価するものであり，ターゲットのラベルが分類器のどの部分を変えるかを測った研究はなく，顔を含む構成はまだ調べられていない．
音声とテキストによるコーパス横断の認識の最近の研究も，コーパス横断のマルチモーダルの手法が少ないと指摘している\cite{ahn2025multitask}．
Gideonら\cite{gideon2021addog}は，いつ汎化し，いつ特化するかの定式化を今後の課題に挙げ，FengとChaspari\cite{feng2023fewshot}は，ソースの性質の影響とソースを組み合わせることの影響を挙げている．
本論文の判断基準と向きの分析と，ラベルの本数に換算したソースの価値は，線形ヘッドについてこれらの問いを扱う．
ラベルを使わない音声感情認識のテスト時適応では，エントロピー最小化と擬似ラベルがほとんどの設定で失敗すると報告されており\cite{dong2026test}，この結果は本論文の比較と一致する．

\emph{少数の標本から推定した線形の規則の誤差．}
標本線形判別の誤分類率の漸近展開は，次元によらない$1/N$の次数の項と，次元に比例する項を含み\cite{okamoto1963asymptotic}，Raudysは後者を実効次元で表した\cite{raudys1998intrinsic}．
本論文の判断基準のコスト$1/(2K)$と向きのコスト$q/K$は，罰則つきのロジスティック回帰についてこの分け方に従い，$q$には測った値を使う．
正解率の式\eqref{eq:twocloud}は，2つのガウス分布のクラスに対する線形の規則の標準的な式である\cite{desilva2023value,dobriban2018highdimensional}．
凍結した特徴と線形ヘッドについては，ドメイン調整回帰が超過リスクをターゲットの平均の誤差と残差で抑え\cite{rosenfeld2022domainadjusted}，分布のずれの下では微調整より線形プロービングが好ましいとされ\cite{kumar2022lpft}，正則化を0に近づける極限ではロジットへの加法的な補正が効果を失う\cite{kini2021labelimbalanced}．
少数のターゲット標本による転移の有限標本の解析は，回帰\cite{ha2026labeled}と1次元の判別\cite{desilva2023value}について存在する．
本論文は，実際のコーパスの対の上の多クラスロジスティック回帰について，対応する量を測る．

\emph{ターゲットの標本数で表したソースの価値．}
事前学習やソースの価値は，ターゲットのデータの実効的な量として文字数で\cite{hernandez2021scaling}，訓練例の数で\cite{lescao2021data}表され，また事前分布の実効標本サイズとして\cite{wiesenfarth2020quantification}，収束の速さの指数として\cite{hanneke2019value}表されてきた．
本論文の式\eqref{eq:value}は，推定した向きのコサインと標本サイズの古典的な関係\cite{raudys1998intrinsic,dobriban2018highdimensional}を逆に解いて，ソースの向きを感情ごとのラベルの本数で表し，本論文はそれをコーパスの対で測る．

\emph{合併と複数のソース．}
複数のドメインからの学習の上界では，ドメイン間の距離は罰則としてだけ働く\cite{bendavid2010theory,blitzer2007learning}．
合併した線形分類器の偏りは，高次元の極限で，タスクごとの切片をもつ中心化したマルチタスクのモデル\cite{tiomoko2021deciphering}と，切片をもたない合併の最小二乗法\cite{yang2025precise}について解かれており，ドメインの情報を明示的に加えると合併を上回りうる\cite{zhu2026domain}．
ニュートン法の1ステップは，点の集合の影響を近似する\cite{pregibon1981logistic,koh2019accuracy}．
これを，ターゲット専用の切片を設けずにターゲットの点を加える場合に当てはめると，式\eqref{eq:lambda}が得られる．
この式は，De Silvaら\cite{desilva2023value}の1次元の判断基準の移動を重みの経路で拡張したものであり，ターゲットがソースから分離できる度合いが，この経路を通じて判断基準を直す．

\emph{特権情報とモダリティ間の蒸留．}
タスクがラベルを共有し，入力のモダリティが異なるとき，追加のモダリティは特権情報（privileged information）である\cite{vapnik2009new,lopezpaz2016unifying}．
モダリティ間の蒸留は，対になったモダリティの間で教師信号を移す\cite{gupta2016cross}．
感情認識では，顔の先生がラベルのない動画で音声の生徒を学習させ\cite{albanie2018emotion}，音声とテキストを融合した先生が音声の生徒を学習させ\cite{srinivasan2022representation}，音声と映像の先生や映像の先生が，ソフトラベルや埋め込みを音声の生徒に渡す\cite{pendyala2025leveraging,hajavi2024audio}．
単一モダリティの先生は融合の生徒を改善しうるし\cite{du2021improving,xue2021multimodal,rakib2025g2d}，融合の出力で単一モダリティのヘッドを学習させることもできる\cite{jaritz2020xmuda}．
これらの研究では，先生をターゲットのコーパスか第3のコーパスで学習するか，生徒が推論時にすべてのモダリティを使う．
本論文の先生はソースコーパスだけで学習した線形分類器であり，学び直さずにターゲットのプールにラベルを付ける．
生徒は，ソースから作った出発点（start）にアンカーした単一モダリティの線形ヘッドである．

\emph{弱いラベル，クラスの釣り合い，アンカー，選択．}
固定した先生が，クラスの異なるタスクをまたぐ表現学習のために\cite{phoo2021selftraining}，また固定した重みによるトランスダクティブな少数ショット分類のために\cite{zanella2024boosting}，ラベルのないデータすべてにソフトラベルを付けた研究があり，モダリティごとの先生をドメインをまたいで単一モダリティに蒸留した研究もある\cite{hatano2024multimodal}．
擬似ラベルをクラス間で釣り合わせることは一般的であるが\cite{huang2022unsupervised,xie2020selftraining,hu2021leveraging,berthelot2022adamatch}，ラベルのない集合のクラスが偏っていると，その釣り合わせが害になりうる\cite{veilleux2021realistic}．
2次の罰則でゼロショット分類器にアンカーした線形プローブ\cite{silvarodriguez2024closer}と仮説転移（hypothesis transfer）\cite{kuzborskij2017fast,li2018explicit}は，重みを事前の解にアンカーし，1つのドメインを除く検証（leave-one-domain-out）はターゲットのラベルなしでハイパーパラメータを選ぶ\cite{gulrajani2021domainbed}．
提案法はこれらの構成要素を組み合わせる．
提案法に固有なのは，ソースだけで学習した別のモダリティの先生，ソースとターゲットが感情の集合を共有するので分類器そのものの学習に使う弱いラベル，重みと切片に別々に置くアンカー，先生と弱いラベルの重みを$K$ごとに選ぶことである．

\emph{弱いラベルによる利得とラベルを使わない選択の理論．}
弱い先生のラベルで学習した生徒の利得は，両者の不一致（misfit）で下から抑えられ\cite{charikar2024quantifying}，両者の表現の食い違いと関係づけられている\cite{dong2025discrepancies,xue2025representations}．
自己訓練は，徐々に進む分布のずれの下で解析されている\cite{kumar2020understanding}．
ガウスの線形モデルでは，最適な計画は観測された応答によらず\cite{chaloner1984optimal,chaloner1995bayesian}，予算が小さいときはランダムな選択や典型的な点の選択を上回るのは難しい\cite{hacohen2022active}．
命題\ref{prop:span}は表現定理\cite{kimeldorf1971results,scholkopf2001generalized}から導かれる．
本論文はこれらの結果を，出発点にアンカーした生徒に用いる．
先生の出発点に対するUARの優位が利得の規模を決め，ラベルを使わない選択は，修正される部分空間を出発点の向きの誤差にそろえられない．

\section{設定と測定の手順}
\label{sec:setting}

本節では，データ，学習器，分割，および第I部と第II部で測る量を定義する．
すべての実験は各反復のテスト話者とターゲットのラベルを共有するので，手法間の差は反復の中で対応がある．

\subsection{コーパスとラベル}
\label{sec:setting-data}

6つのコーパスから7つの群を用いる（表\ref{tab:data}）．
IEMOCAP\cite{busso2008iemocap}の即興のセッションと台本のセッション，CMU-MOSEI\cite{zadeh2018mosei}，MELD\cite{poria2019meld}，AFEW\cite{dhall2012afew}，RAVDESS\cite{livingstone2018ravdess}，CREMA-D\cite{cao2014cremad}である．
IEMOCAPの2種類のセッションを別々の群とするのは，台本のセッションでは10人の話者全員が同じ台詞を読み，テキストの分類器が感情の代わりに台本を認識できてしまうからである．
各発話は，中立（neutral），喜び（happy），悲しみ（sad），怒り（angry）の4つのラベルのいずれかをもつ．
IEMOCAPのhappyとexcitedを統合し，frustrated（IEMOCAP）とcalm（RAVDESS）を除く．
CMU-MOSEIの発話には，6つの強度の評定のうちちょうど1つが$0.5$以上のときにその感情を与え，どれも$0.5$以上でないときに中立を与える．
したがって，CMU-MOSEIの中立のクラスは評定された感情が無いことで定義されるのに対し，IEMOCAPとMELDの中立のクラスは注釈者が割り当てたものである．

音声は7つの群すべてにあり，顔の映像は6つ，書き起こしは4つの群にある．
AFEWには書き起こしが無く，RAVDESSとCREMA-Dでは話者が決まった文を読み上げ，RAVDESSについては顔の特徴を抽出しなかった．
A（音声），V（顔），T（テキスト）を単一モダリティの構成（configuration）と呼び，AV，VT，AT，AVTを融合の構成と呼ぶ．
融合の構成は，そのすべてのモダリティがそろった発話だけを使う．
すべての分割は話者の同一性にもとづいて定める．
CMU-MOSEIには話者の注釈が無い．
チャンネルが分かるのは動画の\NUM{data-mosei-chan-pct}\%であり，それ以外では動画の識別子を話者の代わりに使うので，CMU-MOSEIの話者の分離は不完全である（\ref{sec:limitations-ethics}節）．

\begin{table}[t]
\caption{群，利用できるモダリティ（A：音声，V：顔，T：テキスト），話者のまとまりの数，および除外の後の感情ごとの発話数（音声の特徴での数）．}
\label{tab:data}
\centering
\footnotesize
\setlength{\tabcolsep}{2pt}
\begin{tabular}{@{}llrrrrrr@{}}
\toprule
群 & モダリティ & 話者 & 発話 & 中立 & 喜び & 悲しみ & 怒り \\
\midrule
IEMOCAP 即興 & AVT & \NUM{data-spk-impro} & \NUM{data-n-impro-wavlm} & \NUM{data-impro-neu} & \NUM{data-impro-hap} & \NUM{data-impro-sad} & \NUM{data-impro-ang} \\
IEMOCAP 台本 & AVT & \NUM{data-spk-script} & \NUM{data-n-script-wavlm} & \NUM{data-script-neu} & \NUM{data-script-hap} & \NUM{data-script-sad} & \NUM{data-script-ang} \\
CMU-MOSEI & AVT & \NUM{data-spk-mosei}\textsuperscript{a} & \NUM{data-n-mosei-wavlm} & \NUM{data-mosei-neu} & \NUM{data-mosei-hap} & \NUM{data-mosei-sad} & \NUM{data-mosei-ang} \\
MELD & AVT & \NUM{data-spk-meld} & \NUM{data-n-meld-wavlm} & \NUM{data-meld-neu} & \NUM{data-meld-hap} & \NUM{data-meld-sad} & \NUM{data-meld-ang} \\
AFEW & AV & \NUM{data-spk-afew} & \NUM{data-n-afew-wavlm} & \NUM{data-afew-neu} & \NUM{data-afew-hap} & \NUM{data-afew-sad} & \NUM{data-afew-ang} \\
RAVDESS & A & \NUM{data-spk-ravdess} & \NUM{data-n-ravdess-wavlm} & \NUM{data-ravdess-neu} & \NUM{data-ravdess-hap} & \NUM{data-ravdess-sad} & \NUM{data-ravdess-ang} \\
CREMA-D & AV & \NUM{data-spk-cremad} & \NUM{data-n-cremad-wavlm} & \NUM{data-cremad-neu} & \NUM{data-cremad-hap} & \NUM{data-cremad-sad} & \NUM{data-cremad-ang} \\
\midrule
合計 & & & \NUM{data-n-total-wavlm} & & & & \\
\bottomrule
\end{tabular}

\vspace{2pt}
\raggedright\textsuperscript{a}チャンネルが分かる場合はチャンネル，それ以外は動画の識別子．
\end{table}

\subsection{凍結した特徴と線形分類器}
\label{sec:setting-features}

どのエンコーダも7つの群で微調整しない．
各モダリティの第1のエンコーダで分析を行い，第2のエンコーダで結論が保たれるかを確かめる（\ref{sec:comparison}節）．
音声の特徴はWavLM-Large\cite{chen2022wavlm}（第2：HuBERT-Large\cite{hsu2021hubert}）の最終層を時間で平均し，テキストの特徴はRoBERTa-base\cite{liu2019roberta}（第2：DeBERTa-v3-base\cite{he2023debertav3}）の最終層をトークンで平均する．
顔の特徴は，抜き出した各フレームでYOLOv8n-face\cite{yolov8face}が検出した最大の顔を，VGGFace2\cite{cao2018vggface2}で学習しAffectNet\cite{mollahosseini2019affectnet}で微調整したRyuminaら\cite{ryumina2022facial}のResNet-50（第2：同じ切り出しに対するCLIP ViT-L/14\cite{radford2021clip}の画像エンコーダ）で埋め込み，その埋め込みを発話にわたって平均する．
融合の構成では，各次元を標準化した後に単一モダリティの特徴を連結する．
抽出の詳細は補足資料に示す．

学習器は，訓練データの行で標準化した特徴の上の多項ロジスティック回帰であり，重みに罰則$\|W\|_F^2/(2C)$を課し，切片には罰則を課さない．
クラスの頻度に反比例するクラスの重みにより，学習器は，本論文の主な指標である重みなし平均再現率（unweighted average recall，UAR）を目標とする．
UARのでたらめな予測での水準（chance level）は25\%である．
倍精度のL-BFGS\cite{liu1989lbfgs}で，勾配の絶対値の最大が$10^{-6}$を下回るまで解く．
凍結したエンコーダと線形ヘッドを用いると，判定の規則を分析できる．
感情ごとのラベルが少ないと，微調整は事前学習した特徴を歪めることもあり\cite{kumar2022lpft}，\ref{sec:comparison}節では，本論文のデータで微調整した音声エンコーダが凍結した処理の流れ（パイプライン）を下回ることを示す．

\subsection{分割とラベルの本数}
\label{sec:setting-splits}

各ターゲットの群について，話者を5つの分割（fold）に分け，反復$r$では分割$r \bmod 5$でテストする．
この分割は平均で発話の\NUM{data-test-pct-mean}\%（範囲は\NUM{data-test-pct-lo}--\NUM{data-test-pct-hi}\%）を含む．
残りの発話をプール（pool）とする．
ラベルの本数$K$（label budget）は，プールからランダムに引いたラベル付きのターゲット発話の感情ごとの本数である（$K=25$は$100$本のラベルを意味する）．
引いたラベルは$K$の間で入れ子になっており，どのソースでも同一である．
乱数の種がターゲットと反復だけで決まるからである．
プールの残りはラベルなしで使え，どの手法もテストの発話を学習にも選択にも使わない．
テスト話者はソースからも除く．
これが影響するのはIEMOCAPの2つの群の間だけである．
$K=0$はターゲットのラベルが無いことを表す．

\subsection{3つの学習法と利得の地図}
\label{sec:setting-map}

ターゲットとソースの群のすべての順序対について，各反復で3つの分類器を学習する．
$4K$本のターゲットのラベルで学習する\emph{ターゲットのみ}（target-only），ソースのラベル付き発話すべてで学習する\emph{ソースのみ}（source-only），両方で学習する\emph{合併}（pooled）である．
ソースの\emph{利得}（gain）は，同じ反復における合併のUARからターゲットのみのUARを引いた値であり，テスト話者とラベルを共有する対応のある差である．
\emph{利得の地図}（gain map）は，$K \in \{1, 5, 10, 25, 50\}$とプール全体について，\NUM{map-reps-a}回の反復にわたるUARを，Aでは\NUM{map-pairs-a}，VとAVでは\NUM{map-pairs-v}，テキストを含む各構成では\NUM{map-pairs-t}の順序対について報告する．

ターゲットのみとソースのみの分類器では，話者で分けた3分割の交差検証で$C \in \{0.01, 0.1, 1, 10\}$を選ぶ（ある分割にどれかの感情が欠けるときは$C=1$とする）．
合併の分類器は，同じ反復のソースのみの分類器の$C$を使う．
この$C$の選択にはターゲットのラベルを使わない．
ターゲットのラベルで$C$を検証する方法は，$K$が小さいと失敗する．
感情ごとに1本のラベルでは，ほとんどの反復で検証の点数が同点になり，同点のときの規則が罰則を弱めて合併の分類器の向きを変える．
先行するコーパス横断の研究も，ターゲットに検証に使うラベルが無いときはハイパーパラメータを固定している\cite{gideon2021addog}．
プール全体を使うとき，この規則は合併の分類器を必要以上に制約し，合併の損失を過大に示す．
交差検証で選んだ$C$による地図は補足資料に示す．

\subsection{対ごとの判定の判断基準と向き}
\label{sec:setting-model}

4クラス分類器を，その6つの対ごとの判定を通して分析する．
感情$k$と$l$について，$z = (\mathbf{w}_k - \mathbf{w}_l)^{\top}\mathbf{x} + (b_k - b_l)$をそれらのスコアの差とし，$k$と$l$のターゲット発話の上で，$\mu_k$，$\mu_l$を$z$のクラスごとの平均，$\sigma$を$z$のクラス内の合併標準偏差とする．
\emph{識別力}（separation）$d = (\mu_k - \mu_l)/\sigma$は向き$\mathbf{w}_k - \mathbf{w}_l$だけで決まり，\emph{判断基準のずれ}（threshold offset）$o = (\mu_k + \mu_l)/(2\sigma)$は切片だけで0にできる．
各クラスの中で$z$がガウス分布に従うならば，対の正解率（2つの再現率の平均）は，2つのガウス分布のクラスに対する線形の規則の標準的な式
\begin{equation}
\label{eq:twocloud}
a(d, o) = \tfrac{1}{2}\left[\Phi\!\left(\tfrac{d}{2} + o\right) + \Phi\!\left(\tfrac{d}{2} - o\right)\right]
\end{equation}
で与えられる．
ここで$\Phi$は標準正規分布の分布関数である．
$\Delta$を，ラベル付きのターゲット発話すべての上の線形判別分析で求めた，ターゲットにとって最良の向きに沿った識別力とする．
この基準は分析だけに使う．
この基準に対して，分類器の損失は，判断基準を中点へ動かせば取り戻せる\emph{判断基準による損失}（threshold loss）$a(d, 0) - a(d, o)$と，その後に残る\emph{向きによる損失}（direction loss）$a(\Delta, 0) - a(d, 0)$に分かれる．
両者を6つの対で平均し，対の正解率のポイントで表す．
式\eqref{eq:twocloud}が測った正解率をどれだけよく再現するかは，\ref{sec:part1-what-few-labels-fix}節で確かめる．

\subsection{ハイパーパラメータの選択と報告の仕方}
\label{sec:setting-selection}

第II部の手法とその比較手法（\ref{sec:comparison}節）には，$K$本のターゲットのラベルでは選べないハイパーパラメータがある．
これらは，1つのドメインを除く選択\cite{gulrajani2021domainbed}と同じく，仮のターゲット（pseudo-target）で選ぶ．
ターゲット$T$に対して，別のコーパスに属するソースの群$c$を1つずつ，$T$と$c$のコーパスの外の群から作った手法のターゲットとし，$c$から感情ごとに$K$本のラベルを引く．
選ぶ値は，$c$のテスト話者でのUARを，$c$，最初の5回の反復，各手法が指定するラベルの本数にわたって平均した値を最大にするものである（\ref{sec:part2-mechanism-to-method}節）．
$T$のラベルは，引いた$K$本のほかは選択に入らず，比較手法も同じ方法で調整する．
$T$のテスト話者での最良の値である\emph{オラクル}（oracle）は，選択による損失を測るためだけに使う．

UARの差はパーセントポイントで報告する．
順序対の中（第I部）またはターゲットの中（第II部と\ref{sec:comparison}節）で反復にわたって平均した後，順序対またはターゲットにわたって平均する．
区間は，反復にわたる対応のある差の95\%の$t$区間である．
これらの区間はラベルとテスト話者の引き方による変動を表し，話者やコーパスの再標本化による変動は表さない．
IEMOCAPの10人の話者からは，異なるテストの分割が5つ（即興）または7つ（台本）しかできないので，IEMOCAPのターゲットの反復を独立とみなした区間は狭すぎる．
ソースがIEMOCAPのターゲットに役立つかを判断する箇所では，先に同じ分割を共有する反復を平均する．

\section{第I部：少数のターゲットラベルが最初に直すもの}
\label{sec:part1-what-few-labels-fix}

第I部では，\ref{sec:setting}節の線形ヘッドについて2つの問いを立てる．
感情ごとに$K$本のラベルが線形ヘッドの何を変えるか，その変化のうちどれだけをソースコーパスが代わりに補えるかである．
まず利得の地図を述べ，次に，その下で判断基準と向きの推定のコストが異なる仮定を述べ，最後にそれらのコストから導かれる4つの予言を検証する．

\subsection{利得の地図}
\label{sec:part1-map}

ソースを加えると，ターゲットのラベルが少ないときは役立ち，多いときは害になる（図\ref{fig:gainmap}）．
音声では，$K=1$での平均の利得は$\NUM{map-gain-a-k1}$ポイントであり，\NUM{map-pairs-a}の順序対のうち\NUM{map-gain-pos-a-k1}で95\%区間が0より上にある．
平均の利得は$K=25$までに0を下回り（$\NUM{map-gain-a-k25}$ポイント），プール全体では\NUM{map-gain-neg-a-kall}の順序対で区間が0より下にある．
8つの構成はすべて同じ推移をたどる．
$K=1$での平均の利得は$\NUM{map-gain-v-k1}$（V）から$\NUM{map-gain-avt-k1}$ポイント（AVT）の範囲にあり，$K=25$ではどの構成でも負である．
利得の大きさは，主に2つのコーパスの性質で決まる．
ターゲットごとに1つ，ソースごとに1つの効果をもつ加法モデルは，A，A（HuBERT），V，AVにわたって，利得の分散と，ソースがまだ役立つ最大のラベルの本数の分散の\NUM{map-share-a-lo}--\NUM{map-share-v-hi}\%を説明する．
第I部の残りでは，この推移を，ラベルが分類器の何を変えるかによって説明する．

\begin{figure*}[t]
\centering
\FIG{gainmap}{\textwidth}
\caption{利得の地図．(a)~ターゲットのラベルだけによる学習に対する合併の平均の利得（同じ反復での合併のUARからターゲットのみのUARを引いた値）を順序対にわたって平均し，感情ごとのラベルの本数$K$に対して示す．(b)~利得の95\%区間が0より上にある順序対，0を含む順序対，0より下にある順序対の割合．これらの区間は30回の反復を独立とみなしており，そのためIEMOCAPのターゲットでは狭すぎる．IEMOCAPの2つのセッションの対を含むすべての順序対を用いる．合併の分類器はソースのみの分類器の罰則を用いる．}
\label{fig:gainmap}
\end{figure*}

\subsection{少数のラベルの2つのコスト}
\label{sec:part1-costs}

\ref{sec:setting-model}節の対ごとの判定を，3つの仮定の下で分析する．
\begin{enumerate}
\item[(A1)\hspace{0.5\zw}]各感情の中で，対のスコア差$z$は，2つの感情で共通の分散をもつガウス分布に従う．したがって式\eqref{eq:twocloud}が対の正解率を与える．
\item[(A2)\hspace{0.5\zw}]学習器は重みに有限の$\ell_2$罰則を保ち，切片には罰則を課さない．分離可能なデータで罰則を0に近づけると，ロジットへの加法的な補正は効果を失う\cite{kini2021labelimbalanced}．有限の罰則があれば，切片は自由に動ける．
\item[(A3)\hspace{0.5\zw}]感情ごとの$K$本のラベルは，プールからのランダムな抽出である．
\end{enumerate}

\emph{判断基準．}
向きを固定すると，2つの感情を釣り合わせる切片は，判断基準を$z$の2つのクラス平均の中点に置く．
感情ごとに$K$本のラベルから推定すると，この中点の分散は$\sigma^2/(2K)$なので，残るずれは$\mathbb{E}[o^2] = 1/(2K)$を満たす．
式\eqref{eq:twocloud}を$o$について2次まで展開すると，残る判断基準による損失の期待値
\begin{equation}
\label{eq:thrcost}
\mathbb{E}\left[a(d, 0) - a(d, o)\right] \approx \frac{d}{4}\,\phi\!\left(\frac{d}{2}\right)\frac{1}{2K} \le \frac{\phi(1)}{4K}
\end{equation}
が得られる．
ここで$\phi$は標準正規分布の密度関数であり，$d$についての最大は$d=2$で達成される．
したがって少数のラベルは，判断基準のずれがどれほど大きくても，その大部分を取り除く．

\emph{向き．}
向きを推定するにはベクトルを推定する必要があり，その誤差は実効次元$q$と$K$の比に応じて縮む．
感情ごとに$K$本のラベルから推定した線形の規則では，識別力は
\begin{equation}
\label{eq:dircost}
d(K) = \frac{\Delta}{\sqrt{1 + q/K}}
\end{equation}
のようにふるまう．
これは，線形判別の誤分類率の漸近展開が向きの項でとる形である\cite{okamoto1963asymptotic,raudys1998intrinsic}．
$q$は，ターゲットのみの分類器の識別力をラベルの本数に沿って並べたものに式\eqref{eq:dircost}を当てはめて測る．
当てはめた$q$は，WavLMの音声で$\NUM{p1-qdir-a}$，HuBERTの音声で$\NUM{p1-qdir-ah}$，AVで$\NUM{p1-qdir-av}$，顔で$\NUM{p1-qdir-v}$である．
したがって，音声でターゲットにとって最良の向きに匹敵する向きを得るには，感情ごとに20本程度のラベルが要る．
このように，本論文で調べた$K$の範囲では，ラベルは判断基準を少ない費用で直し，音声の向きを遅く直す．
ここから4つの予言が導かれる．
最初の2つを検証する構成ごとの測定値を表\ref{tab:part1}にまとめる．
(P1)~ソースのみの分類器の損失のうち，少数のラベルで取り戻す部分の大半は判断基準による損失である．
(P2)~利得のうち判断基準から生じる割合は，$K$が増えると下がる．
(P3)~合併は，ターゲットが受ける重みに比例して判断基準を直す（\ref{sec:part1-pooling}節）．
(P4)~ソースは，$q$と，ソースの向きがターゲットにどれだけ合うかで決まる本数のターゲットのラベルの代わりになる（\ref{sec:part1-value}節）．

\begin{figure}[t]
\centering
\FIG{twocloud}{\columnwidth}
\caption{対ごとの判定の判断基準と向き．(a)~スコア差$z$に沿った2つのガウス分布のクラス．判断基準のずれ$o$は判定の点を中点から遠ざけ，識別力$d$は向きで決まる．(b)~ターゲットにとって最良の線形の規則に対するソースのみの分類器の損失を，判断基準による損失と向きによる損失に分けて，対の正解率のポイントで示す．6つの感情の対と順序対にわたって平均した．}
\label{fig:twocloud}
\end{figure}

\subsection{ラベルが直すもの}
\label{sec:part1-fix}

2つのガウス分布のモデル（two-cloud model）は，本論文で学習した分類器の対の正解率を再現する．
すべての感情の対と順序対にわたって，測った$d$と$o$から予言した正解率と測った正解率の相関は\NUM{p1-e13-fit-r-lo}--\NUM{p1-e13-fit-r-hi}であり，平均絶対誤差は\NUM{p1-e13-fit-mae-lo}--\NUM{p1-e13-fit-mae-hi}ポイントである．

\emph{顔を除いて，ソースのみの分類器は主に向きによって損をする．}
WavLMの音声では，判断基準による損失は\NUM{p1-e13-loss-offset-a}ポイント，向きによる損失は\NUM{p1-e13-loss-dir-a}ポイントである（図\ref{fig:twocloud}b）．
顔ではそれぞれ\NUM{p1-e13-loss-offset-v}と\NUM{p1-e13-loss-dir-v}ポイント，テキストでは\NUM{p1-e13-loss-offset-t}と\NUM{p1-e13-loss-dir-t}ポイントである．
判断基準による損失は，構成にわたって全損失の\NUM{p1-e29-lossshare-lo}--\NUM{p1-e29-lossshare-hi}にすぎない．

\emph{判断基準による損失は，大部分がコーパスの平均のずれである．}
中心化（centering）は各コーパスのラベルのない発話の平均を引く操作であり，ラベルを使わない．
中心化は，ソースのみのUARを，WavLMの音声で\NUM{p1-center-srconly-a}ポイント，HuBERTの音声で\NUM{p1-center-srconly-ah}ポイント，顔で\NUM{p1-center-srconly-v}ポイント上げ，音声の\NUM{p1-center-srconly-pairs-a-n}の対のうち\NUM{p1-center-srconly-up-a-n}で上げる（3回の反復）．
重みを固定し，ソースのみの分類器の切片だけを感情ごとに$K=5$本のラベルで当てはめ直すと，WavLMの音声で\NUM{p1-d1-gain-a-k5}ポイントの利得があり，これは同じラベルで合併が得る\NUM{p1-merged-gain-a-k5}ポイントに近い．
中心化の後では，この2つの利得は\NUM{p1-center-d1-gain-a-k5}と\NUM{p1-center-merged-gain-a-k5}ポイントに縮む（補足資料の図S2）．

\emph{ラベルはまず判断基準に使われる（P1，P2）．}
ソースのみの分類器に対する合併の分類器の改善を，測った$d$と$o$から式\eqref{eq:twocloud}で計算した判断基準の部分と向きの部分の和として書くと，単一モダリティの音声と顔では$K \le 5$で判断基準の部分が大半を占める．
その割合は，WavLMの音声，HuBERTの音声，顔の$K=1$と$K=5$にわたって\NUM{p1-e29-critshare-kle5-single-lo}--\NUM{p1-e29-critshare-kle5-single-hi}である（図\ref{fig:thresholdshare}aと表\ref{tab:part1}．この割合はIEMOCAPの2つのセッションの対を除いて求めた）．
この割合は損失に占める判断基準の割合よりずっと大きいので，ラベルは誤差のうち少ない費用で直せる部分を先に直す．
予言どおり，この割合は$K$とともに下がり，$K=10$ではHuBERTの音声で\NUM{p1-e29-critshare-k10-ah}である．
融合の構成とテキストの構成では判断基準による損失が小さく，$K=5$では判断基準の部分は利得のおよそ半分かそれ以下である（\NUM{p1-e29-critshare-k5-fused-text-lo}--\NUM{p1-e29-critshare-k5-fused-text-hi}，AVでは\NUM{p1-e29-critshare-k5-av}）．
$K=1$では，AV（\NUM{p1-e29-critshare-k1-av}）とVT（\NUM{p1-e29-critshare-k1-vt}）で判断基準の部分がまだ半分を超える．
同じ特徴の上の2層パーセプトロンでは，$K=5$での判断基準の割合は音声で\NUM{p1-e30-critshare-k5-audio-lo}--\NUM{p1-e30-critshare-k5-audio-hi}，顔で\NUM{p1-e30-critshare-k5-v-wd2}--\NUM{p1-e30-critshare-k5-v-wd4}である．
したがってこの傾向は線形ヘッドに固有ではない．
ただし，HuBERTの音声では割合が半分を下回る．

\emph{向きはソースが決めたところにとどまる．}
式\eqref{eq:dircost}の下では，ソースの向きと感情ごとに10本のラベルを最適に組み合わせても，音声の向きによる損失のうち取り戻せるのは\NUM{p1-dirloss-recov-bound-k10-audio-lo}--\NUM{p1-dirloss-recov-bound-k10-audio-hi}\%である．
ソースの分類器から出発してラベルで学び直す仮説転移が取り戻すのは\NUM{p1-dirloss-recov-ht-k10-audio-lo}--\NUM{p1-dirloss-recov-ht-k10-audio-hi}\%である．

\begin{figure}[t]
\centering
\FIG{thresholdshare}{\columnwidth}
\caption{ラベルが何に使われるか．(a)~ソースのみの分類器に対する合併の分類器の改善のうち判断基準から生じる割合を，$K$に対して示す（IEMOCAPの2つのセッションの対を除いた順序対の中央値．10回の反復）．実線は単一モダリティの音声と顔，破線は融合の構成とテキストの構成，黒い水平線は2分の1を示す．(b)~感情ごとに$K$本のラベルから推定した後に残る誤差の期待値．判断基準のずれは$1/(2K)$で，向きの誤差は$q/K$で下がる．$q$は構成ごとに当てはめた値である．}
\label{fig:thresholdshare}
\end{figure}

\begin{table}[t]
\caption{構成ごとの第I部の量．ソースのみの分類器の損失は対の正解率のポイントで示す．$q$は式\eqref{eq:dircost}の向きの実効次元である．最後の3列は，合併による改善のうち判断基準から生じる割合である（順序対の中央値）．}
\label{tab:part1}
\centering
\footnotesize
\setlength{\tabcolsep}{3.5pt}
\begin{tabular}{@{}lrrrrrr@{}}
\toprule
 & \multicolumn{2}{c}{損失（ポイント）} & & \multicolumn{3}{c}{判断基準の割合} \\
\cmidrule(lr){2-3}\cmidrule(l){5-7}
構成 & 判断基準 & 向き & $q$ & $K{=}1$ & $K{=}5$ & $K{=}10$ \\
\midrule
A & \NUM{p1-e13-loss-offset-a} & \NUM{p1-e13-loss-dir-a} & \NUM{p1-qdir-a} & \NUM{p1-e29-critshare-k1-a} & \NUM{p1-e29-critshare-k5-a} & \NUM{p1-e29-critshare-k10-a} \\
A（HuBERT） & \NUM{p1-e13-loss-offset-ah} & \NUM{p1-e13-loss-dir-ah} & \NUM{p1-qdir-ah} & \NUM{p1-e29-critshare-k1-ah} & \NUM{p1-e29-critshare-k5-ah} & \NUM{p1-e29-critshare-k10-ah} \\
V & \NUM{p1-e13-loss-offset-v} & \NUM{p1-e13-loss-dir-v} & \NUM{p1-qdir-v} & \NUM{p1-e29-critshare-k1-v} & \NUM{p1-e29-critshare-k5-v} & \NUM{p1-e29-critshare-k10-v} \\
AV & \NUM{p1-e13-loss-offset-av} & \NUM{p1-e13-loss-dir-av} & \NUM{p1-qdir-av} & \NUM{p1-e29-critshare-k1-av} & \NUM{p1-e29-critshare-k5-av} & \NUM{p1-e29-critshare-k10-av} \\
T & \NUM{p1-e13-loss-offset-t} & \NUM{p1-e13-loss-dir-t} & -- & \NUM{p1-e29-critshare-k1-t} & \NUM{p1-e29-critshare-k5-t} & \NUM{p1-e29-critshare-k10-t} \\
VT & \NUM{p1-e13-loss-offset-vt} & \NUM{p1-e13-loss-dir-vt} & -- & \NUM{p1-e29-critshare-k1-vt} & \NUM{p1-e29-critshare-k5-vt} & \NUM{p1-e29-critshare-k10-vt} \\
AT & \NUM{p1-e13-loss-offset-at} & \NUM{p1-e13-loss-dir-at} & -- & \NUM{p1-e29-critshare-k1-at} & \NUM{p1-e29-critshare-k5-at} & \NUM{p1-e29-critshare-k10-at} \\
AVT & \NUM{p1-e13-loss-offset-avt} & \NUM{p1-e13-loss-dir-avt} & -- & \NUM{p1-e29-critshare-k1-avt} & \NUM{p1-e29-critshare-k5-avt} & \NUM{p1-e29-critshare-k10-avt} \\
\bottomrule
\end{tabular}
\end{table}

\subsection{合併が判断基準を直す仕組み}
\label{sec:part1-pooling}

合併は，ターゲット用の別の切片をもたずに判断基準を直す．
その仕組みは，ソースのみの解からのニュートン法の1ステップで示せる（P3）．
$N_s$と$N_t$をソースとターゲットの訓練発話の数，$\kappa_s$と$\kappa_t$をソースのみの解における損失の平均の曲率，$A_s$をソースの特徴の曲率で重みづけした散布行列，$\boldsymbol{\delta}$をターゲットの平均と曲率で重みづけしたソースの平均の差とする．
1ステップの後には，ソースのみのずれ$o_s$のうち割合$1 - \lambda$が残る．
ただし
\begin{equation}
\label{eq:lambda}
\lambda = \frac{\rho}{1 + \rho},\quad \rho = N_t\kappa_t\left[\frac{1}{N_s\kappa_s} + \boldsymbol{\delta}^{\top}\left(A_s + \frac{I}{C}\right)^{-1}\boldsymbol{\delta}\right]
\end{equation}
である．
括弧の中の第1項は切片の経路であり，ターゲットはこの経路を通じて，訓練の重みに占める自身の割合に比例して判断基準を直す．
第2項は重みの経路である．
ソースがほとんど変動しない向きに沿ってターゲットがソースから線形に分離できるとき，合併の分類器はその向きに重みを置き，その重みはターゲットの発話にだけ加わる切片として働く．
この形は，点の集合の影響（group influence）\cite{koh2019accuracy,pregibon1981logistic}を加える点について書いたものであり，その切片の経路は，$n$個のターゲットの点と$m$個のソースの点を合併した1次元の判別の判断基準の移動$m/(n+m)$に帰着する\cite{desilva2023value}．

合併の分類器のターゲット上の$d$と$o$から計算した，測った修正$\lambda = 1 - o_m/o_s$は，式\eqref{eq:lambda}が予言する順に並ぶ．
ただし，この式は修正の平均を過大に見積もる（補足資料の図S1）．
順序対と$K \in \{1, 5, 10\}$にわたって，予言した$\lambda$と測った$\lambda$の順位相関は，A，A（HuBERT），V，AVにわたって\NUM{p1-e15-lam-rho-lo}--\NUM{p1-e15-lam-rho-hi}である（10回の反復）．
完全なニュートン法の1ステップは修正を順位相関\NUM{p1-e16-lam-rho-lo}--\NUM{p1-e16-lam-rho-hi}で並べ，VとAVでは向きの変化も予言する（\NUM{p1-e16-dir-rho-video-lo}--\NUM{p1-e16-dir-rho-video-hi}）．
2次の展開の精度が低い音声では，向きの変化の予言はそれほど当たらない（\NUM{p1-e16-dir-rho-audio-lo}--\NUM{p1-e16-dir-rho-audio-hi}）．
重みの経路は，コーパス間の距離が利得を予言しない理由を説明する．
距離が大きいとソースのみの正解率は下がるが，合併が判断基準を直す経路も広がる．
式\eqref{eq:lambda}の導出と，それを確かめるシミュレーションは補足資料に示す．

\subsection{ソースの価値}
\label{sec:part1-value}

ソースは向きをもたらし，その価値はターゲットのラベルの本数で表せる（P4）．
$e$を，ターゲットにとってのソースの向きの効率，すなわちソースの向きがターゲットの最良の識別力に掛かる係数（$d = e\Delta$）とする．
式\eqref{eq:dircost}の下では，$K$本のラベルから推定した向きの雑音と信号の比は$q/K$であり，ソースの向きの比は$(1 - e^2)/e^2$である．
2つの独立な推定を最適な重みで組み合わせると標本サイズが足し合わされる．
したがって，ソースから出発して$K$本のラベルで学び直すことは，ターゲットだけで感情ごとに$K + n$本のラベルで学習することと同じようにふるまう．
ここで
\begin{equation}
\label{eq:value}
n = q\,\frac{e^2}{1 - e^2}
\end{equation}
である．
式\eqref{eq:value}は，推定した向きのコサインと標本サイズの古典的な関係\cite{raudys1998intrinsic,dobriban2018highdimensional}を逆に解いたものである．
逆に解くことでソースの向きがラベルの本数に換算され，$n$は$K$によらない．

$n$は，仮説転移のUARがターゲットのみの学習曲線のどこに当たるかを求め，そこから$K$を引いて測る（図\ref{fig:sourcevalue}）．
音声のソース1つでは，$n$は感情ごとに\NUM{p1-e32-n1src-audio-lo}--\NUM{p1-e32-n1src-audio-hi}本であり，$K=5$から$K=25$までほぼ一定である．
顔では\NUM{p1-e32-n1src-v-k5}--\NUM{p1-e32-n1src-v-k10}本である．
すべてのソースを合併したものから出発すると，WavLMの音声で，$n$は$K=5$と$K=10$でそれぞれ感情ごとに\NUM{p1-e32-nall-a-k5}本と\NUM{p1-e32-nall-a-k10}本である．
音声の対にわたって，予言した$n$と測った$n$の順位相関は\NUM{p1-e32-rho-audio-lo}--\NUM{p1-e32-rho-audio-hi}であり，測った値は予言した値の\NUM{p1-e32-ratio-audio-lo}--\NUM{p1-e32-ratio-audio-hi}倍である．
すべてのソースを合併した場合，この比は\NUM{p1-e32-allsrc-ratio-lo}--\NUM{p1-e32-allsrc-ratio-hi}である．

\begin{figure}[t]
\centering
\FIG{sourcevalue}{\columnwidth}
\caption{ソースの価値．予言した値$n = q e^2/(1-e^2)$と，ターゲットのみの学習曲線の上で測った値を，感情ごとのラベルの本数で，単一のソースとすべてのソースを合併したものについて比べる（$K=5$と$K=10$．IEMOCAPの2つのセッションの対を除く）．対角線は一致を表す．学習曲線の範囲を超える点は上の帯に，$n \le 0$の点は下の帯に示す．}
\label{fig:sourcevalue}
\end{figure}

\subsection{第I部のまとめ}
\label{sec:part1-summary}

単一モダリティの音声と顔の特徴では，少数のターゲットのラベルが判断基準を直し，中心化はラベルを使わずに判断基準の誤差の大部分を取り除く．
費用がかかるのは向きである．
音声でターゲットの最良の向きに近づくには感情ごとに$q$本程度のラベルが要り，ソース1つはそのうちおよそ10本分を補い，すべてのソースを合わせるとその2倍から3倍を補う．
したがって，$K \le 5$でさらに改善するには，ラベルとソースがこの形で与えるよりも多くの向きの情報が要る．
第II部では，その情報をどこから得られるかを導く．

\section{第II部：仕組みから手法へ}
\label{sec:part2-mechanism-to-method}

第I部の結果，状況は次のとおりである．
中心化と切片の当てはめ直しが判断基準を扱い，ソースは向きについておよそ10本のラベル分を補い，残る向きの誤差が主な損失である．
第II部では，ソースから作れる最も強い出発点を構成し，少数のラベルも，どの発話にラベルを付けるかをラベルを使わずに選ぶことも，その向きをあまり動かせないことを示す．
そのうえで，別のモダリティを見る先生がラベルのないプール全体に付ける弱いラベルによって，足りない向きの情報を補う（図\ref{fig:method}）．

\begin{figure*}[t]
\centering
\FIG{method}{\textwidth}
\caption{提案法．ソースとターゲットの特徴は別々に中心化する．ターゲットのコーパスの外のすべてのソースから出発点の分類器$(W_0, b_0)$を作り，生徒はそこから出発し，重みと切片に別々の$\ell_2$アンカーをもつ．基本の手法では，ソースだけで学習した先生（単一モダリティの生徒に対しては，生徒自身のモダリティと生徒がもたないモダリティを含む，その構成の融合の出発点）が，ラベルのないターゲットの発話すべてに，感情ごとの合計をそろえたソフトラベルを付ける．この弱いラベルは，$4K$本の本当のラベルと並んで，合計の重み$m$で生徒の損失に入る．出発点の5つの値（アンカー$C'$と$C_b$を含む）は仮のターゲットで1度だけ選び，先生と$m$は$K$ごとに選ぶ．拡張版は，先生の候補として，$K$本のラベルで学び直した融合の生徒と，プールで学習した単一モダリティの生徒の積を加え，弱いラベルの重みづけ（釣り合わせる，釣り合わせない，ハード）を$m$とともに$K$ごとに選ぶ．}
\label{fig:method}
\end{figure*}

\subsection{出発点：中心化とアンカーつきの転移}
\label{sec:part2-start}

生徒は中心化した特徴を使う．
各ソースはそれ自身の平均で，ターゲットはそのラベルのないプールの平均で中心化し，これによって判断基準による損失の大部分が取り除かれる（\ref{sec:part1-fix}節）．
出発点$(W_0, b_0)$は，ターゲットのコーパスの外のすべてのソースの向きを合わせたものである．
各ソース$s$は，閉じた形の判別の向き$W_s(\rho_s) = ((1-\rho_s)\Sigma_s + \rho_s \tau_s I)^{-1}\tilde{M}_s$を与える．
ここで$\tilde{M}_s$はその中心化したクラス平均を並べた行列，$\Sigma_s$はその共分散行列，$\tau_s$はその固有値の平均である．
これらの向きを$N_s^{\gamma}$に比例する重みで足し合わせ，切片はクラスの事前確率が等しいと仮定して決める．
生徒は次の問題を解く．
\begin{equation}
\label{eq:student}
\min_{W, b}\ \sum_{i \in \mathcal{L}} \ell\!\left(W \mathbf{z}_i + b, y_i\right) + \frac{\|W - W_0\|_{M}^2}{2C'} + \frac{\|b - b_0\|^2}{2C_b}
\end{equation}
ここで$\mathcal{L}$は$4K$本のラベル付き発話の集合であり，ノルム$\|\cdot\|_M$は，ラベルのないターゲットのプールの散布行列から混合の重み$\rho_M$で作った計量$M$を用いる．
重みと切片に別々のアンカーを置くのは第I部の結果による．
少数のラベルで切片はよく決まるので切片には弱いアンカーで足り，少数のラベルでは重みがうまく決まらないので重みには強いアンカーが要る．
5つの値$(\gamma, \rho_s, C_b, C', \rho_M)$は，それぞれ3つの候補から仮のターゲットで選ぶ（\ref{sec:setting-selection}節）．

この構成の各段階はUARを上げる．
中心化を伴う，合併したソースからの仮説転移は，音声の$K \in \{1, 5, 10\}$で，ターゲットのみの学習，合併，中心化した合併，ターゲットの重みを選んだ合併のうち最良のものを\NUM{p2-e31-p4-vs-best-audio-lo}--\NUM{p2-e31-p4-vs-best-audio-hi}ポイント上回る（30回の反復）．
5つの値を選ぶと，$K=5$でA，A（HuBERT），V，AVにわたって\NUM{p2-e39-vs-p4-k5-lo}--\NUM{p2-e39-vs-p4-k5-hi}ポイント上がり，ソースを等しく重みづけした出発点に対しては\NUM{p2-e39-vs-d1-k5-lo}--\NUM{p2-e39-vs-d1-k5-hi}ポイント上がる（30回の反復）．
この結果を，5つの値を選んだ出発点（five-value start）と呼ぶ．
第II部と\ref{sec:comparison}節では，すべての追加をこの出発点を基準に測る．

\subsection{向きの情報を得られないところ}
\label{sec:part2-limits}

$W^{*}$をターゲットのプールのすべてのラベルを使ったときに得られる生徒とし，$E_0 = W^{*} - W_0$を出発点の向きの誤差とする．
$K$本のラベルが$E_0$に対してできることは，2つの事実によって制限される．

\begin{proposition}[張る空間]
\label{prop:span}
式\eqref{eq:student}の解の$W - W_0$の各行は，$4K$個のラベル付きの点$M^{-1}\mathbf{z}_i$（$i \in \mathcal{L}$）が張る空間に含まれる．
\end{proposition}

損失は積$W\mathbf{z}_i$を通してだけ$W$に依存するので，この命題は，$W - W_0$に表現定理\cite{kimeldorf1971results,scholkopf2001generalized}を当てはめた直接の帰結である．
修正できるのは，$E_0$のうち高々$4K$次元の部分空間への射影だけであり，$E_0$のどの部分がその部分空間に入るかは，どの発話がラベルをもつかによる．

\begin{proposition}[ラベルを使わない計画]
\label{prop:design}
損失を$W_0$のまわりのガウスの尤度で，アンカーを事前分布$W \sim \mathcal{N}(W_0, C' M^{-1})$で近似する．
この事前分布の下で$E_0$に特に好まれる向きが無いならば，任意の$4K$個の点が張る空間が捉える$\|E_0\|^2$の割合の期待値は，点のラベルを使わずに選んだどの選び方でも同じである．
\end{proposition}

ガウスの線形モデルでは事後共分散が観測されたラベルによらないので，最良の計画もラベルによらない\cite{chaloner1984optimal,chaloner1995bayesian}．
ラベルを使わない選び方でも，たとえば点を散らすことで事後分散を減らすことはできるが，$E_0$の向きを知らなければ，修正される部分空間を$E_0$にそろえることはできない．
命題\ref{prop:design}はこの近似に依拠するので，直接検証する．

感情ごとに$K=5$本のラベルが張る空間が捉える$\|E_0\|^2$の割合を，ラベルを使わない7つの選択規則と，すべての候補のラベルを知る教示の規則（teaching rule）の下で測った．
7つの規則は，別のランダムな順序，分類器のマージン，予測クラスの上のherding，k-centerのコアセット，典型性によるクラスタリング，白色化した座標での2つの計画である．
ランダムな選択に対する比（捉える割合をそれぞれターゲットで平均してから取った比）で測ると，ラベルを使わない規則が捉える量は\NUM{p2-e58-ratio-unlab-lo}--\NUM{p2-e58-ratio-unlab-hi}倍，教示の規則が捉える量は\NUM{p2-e58-ratio-lll-lo}--\NUM{p2-e58-ratio-lll-hi}倍である（図\ref{fig:capture}）．
ターゲットのラベルを使わずに，ソース間の食い違いや仮のターゲットで学んだ修正から作った部分空間は，$\|E_0\|^2$の\NUM{p2-e58-r-uaub-lo}--\NUM{p2-e58-r-uaub-hi}\%を捉える．
選択の結果もこれと一致する．
5つの値を選んだ生徒では，標準的な能動学習の規則を含め，ラベルを使わないどの規則も$K=5$でランダムな選択を1ポイント上回らない（最大の差は\NUM{p2-cond2-e39-maxgain-k5}）．
一方，より弱い生徒を5つの値を選んだ生徒に置き換えると，ランダムな選択そのものが\NUM{p2-cond2-student-random-k5-lo}--\NUM{p2-cond2-student-random-k5-hi}ポイント上がる．
したがって，よい出発点に残る向きの誤差はターゲットに固有であり，ラベルを使わない選択からは見えない．
その修正にはラベルの情報が必要であり，しかも$4K$本のラベルが与えるよりずっと多くの情報が必要である．

\begin{figure}[t]
\centering
\FIG{capture}{\columnwidth}
\caption{ラベルを使わない選択は向きの誤差を見つけない．感情ごとに$K=5$本のラベルが張る空間が捉える出発点の向きの誤差$E_0$の割合を，ランダムな選択に対する比で，ラベルを使わない7つの規則と，すべての候補のラベルを知る教示の規則（上側の参照）について，A，V，Tで示す．点はターゲットの平均どうしの比であり，ひげはターゲットごとの比の範囲を示す．10回の反復．}
\label{fig:capture}
\end{figure}

\subsection{プール全体への弱いラベル}
\label{sec:part2-weak}

プールのすべての発話を予測する先生は，数千の発話についてラベルの情報を与える．
$Q \in [0,1]^{N_u \times 4}$を，プールの$N_u$個のラベルのない発話に対する先生のクラス確率を並べた行列とする．
$Q$の各列を，どの感情も同じ合計の重みを受けるように拡大縮小し，行列全体を合計の重みが$m$になるように拡大縮小する．
生徒は，式\eqref{eq:student}にプールにわたる$\sum_{i}\sum_{c} Q_{ic}\,\ell(W\mathbf{z}_i + b, c)$を加えたものを最小化する．
先生の誤りはプールにわたって平均されて打ち消し合い，生徒を動かすのは，先生と出発点の判定が食い違う発話の集合である．
先生と$m$は，弱いラベルを使わない選択肢を含む小さな格子から，$K$ごとに仮のターゲットで選ぶ．

先生の有用さには自然な尺度がある．
$J$を，プールでの先生のUARから出発点のUARを引いた値とする．
定義により，$J$は，先生が正しく出発点が誤って分類する発話の割合から，その逆が成り立つ発話の割合を引いたものを，クラスで平均した値に等しい．
出発点にアンカーしたまま先生のラベルに合わせる生徒は，先生の優位をせいぜい近似的に再現できるだけである．そこで本論文は，定理によらない見込みとして，生徒の利得のうち先生の情報から来る部分は$J$を下回ると考える．
列の正規化は，ラベルを使わない第2の効果を加える．
出発点自身の予測を感情の間で釣り合わせ直したものが，それだけで有用な先生として働く．

同じモダリティの先生がもたらすものは少ない．
ソースごとの分類器の投票を，仮のターゲットでのそれぞれの混同の傾向で重みづけして組み合わせると信念（belief）が得られ，5つの値を選んだ出発点に対するその$J$は，WavLMの音声で\NUM{p2-e52-q-minus-p-a}ポイントである．
この信念からの弱いラベルは，5つの値を選んだ出発点を，音声では$K = 0, 1, 5$で\NUM{p2-e59k-a-k0}，\NUM{p2-e59k-a-k1}，\NUM{p2-e59k-a-k5}ポイント，顔では\NUM{p2-e59k-v-k0}，\NUM{p2-e59k-v-k1}，\NUM{p2-e59k-v-k5}ポイント上げ（10回の反復），利得は$K=25$までに消える．
音声の$K=0$では，出発点自身の予測の釣り合わせ直しが\NUM{p2-e59k-ctrl-a-k0}ポイントを占め，信念の情報は$K \le 5$にわたって\NUM{p2-e59k-info-a-lo}--\NUM{p2-e59k-info-a-hi}ポイントを加える．
これは信念の$J$を下回る．
テキストでは信念は出発点よりよくなく，利得はほぼ0である．

\subsection{別のモダリティの先生}
\label{sec:part2-crossmodal}

ターゲットの発話が他のモダリティも伴うとき，それらのモダリティを見る，ソースで学習した融合の分類器は，別の向きを知っている．
これは，タスクがラベルを共有し入力のモダリティが異なる，特権情報の設定である\cite{lopezpaz2016unifying,vapnik2009new}．
単一モダリティの生徒に対しては，その構成の融合の出発点を先生とする．
この先生は\ref{sec:part2-start}節の手順でソースだけから学習し，ターゲットのデータで学び直さない．
配備した生徒は自分のモダリティだけを使う．
この先生の$J$は大きい．
AVの構成での2回の反復による予備的な確認では，プールで融合の出発点は音声の出発点をUARで\NUM{p2-e60pre-juar-av-a-fused}ポイント上回り，モダリティごとの予測の他の組み合わせは音声の出発点を\NUM{p2-e60pre-juar-av-a-lo}--\NUM{p2-e60pre-juar-av-a-hi}ポイント上回る．
一方，顔の出発点だけでは音声の出発点を\NUM{p2-e60pre-juar-av-a-v}ポイント下回る．
融合の生徒に対しては，単一モダリティの予測の組み合わせは融合の出発点とほとんど変わらない（同じ確認で，AVでの$J$は\NUM{p2-e60pre-juar-av-f-poe}から\NUM{p2-e60pre-juar-av-f-pool}ポイント）．
そこで2段階の参照（two-stage reference）を用いる．
まず単一モダリティの生徒をターゲットのプールで融合の先生から学習し，次にそれらの予測の積で融合の生徒を教える．

$K \le 5$では，別のモダリティからの弱いラベルは，すべての生徒を5つの値を選んだ出発点より上に上げる（図\ref{fig:ladder}に音声の生徒を示す）．
$K=0$での単一モダリティの生徒の利得は，AVで\NUM{p2-e60-av-a-k0}（音声）と\NUM{p2-e60-av-v-k0}ポイント（顔），VTで\NUM{p2-e60-vt-t-k0}（テキスト），ATで\NUM{p2-e60-at-t-k0}（テキスト），AVTで\NUM{p2-e60-avt-a-k0}（音声），\NUM{p2-e60-avt-t-k0}（テキスト），\NUM{p2-e60-avt-v-k0}ポイント（顔）である．
\NUM{p2-e60-m1-n}の単一モダリティの生徒はすべて，$K=0$と$K=5$の両方で利得を得る．
融合の生徒は$K=0$で\NUM{p2-e60-f-k0-lo}--\NUM{p2-e60-f-k0-hi}ポイントの利得を得る．
音声とテキストの生徒では，これらの利得は同じモダリティの先生による利得を上回る．
顔の生徒では小さい（$K=0$で，VTとAVTでの\NUM{p2-e60-vt-v-k0}ポイントと\NUM{p2-e60-avt-v-k0}ポイントに対し，同じモダリティの先生では\NUM{p2-e59k-v-k0}ポイント）．
$J$を測ったAVでは，前の小節の尺度が見込むとおり，音声の生徒の利得はその先生の$J$を下回る．
ラベルが増えると利得は縮み，$K \ge 10$では\NUM{p2-e60-m3-students-n}の生徒のうち\NUM{p2-e60-m3-loss-n}が5つの値を選んだ出発点を下回る（最も悪い差は\NUM{p2-e60-m3-worst}ポイント）．
拡張した候補群（extended family）は，先生の候補として$K$本のラベルで学び直した融合の生徒と2段階の参照を加え，弱いラベルの重みづけ（釣り合わせる，釣り合わせない，ハード）を選択の対象に加えたものである．
拡張した候補群は，比較で評価するAVTの音声の生徒とAVの融合の生徒について，これらの損失をなくす（\ref{sec:comparison}節）．
VTのテキストの生徒とATの音声の生徒は，拡張した候補群では実行していない．

\begin{figure}[t]
\centering
\FIG{ladder}{\columnwidth}
\caption{音声の生徒における提案法の段階ごとの結果．(a)~ターゲットのみの学習，合併（単一のソースにわたる平均），合併したソースからの仮説転移（HT），5つの値を選んだ出発点のUARを$K$に対して示す（7つのターゲット，30回の反復）．(b)~同じ実験の5つの値を選んだ出発点からの差を，同じモダリティの信念からの弱いラベル（7つのターゲット），AVの融合の先生（6つのターゲット）とAVTの融合の先生（4つのターゲット）からの弱いラベル，AVTでの拡張した候補群（4つのターゲット）について，10回の反復で示す．融合の構成はそのすべてのモダリティがそろった発話だけを残すので，出発点は構成ごとに異なる．}
\label{fig:ladder}
\end{figure}

\section{標準的な手法との比較}
\label{sec:comparison}

\subsection{比較手法と手順}
\label{sec:comparison-protocol}

9つの比較手法（competitor）は，Lee\cite{lee2013pseudolabel}の形の，生徒自身の確信度の高い予測による自己訓練，確信度がしきい値を超える先生のハードな擬似ラベル，エントロピー最小化\cite{grandvalet2004semisupervised}，温度つきの蒸留\cite{hinton2015distilling}，一般化蒸留\cite{lopezpaz2016unifying}，共訓練\cite{blum1998combining}，共正則化\cite{sindhwani2005coregularization}，期待値最大化（EM）によるクラスの事前確率の調整\cite{saerens2002adjusting}，出発点の出力の再均衡化である．
分類器を学習し直す手法はどれも，同じ5つの値を選んだ出発点から，同じ$K$本のラベルと同じラベルのないプールで，同じアンカーつきの線形の生徒（式\eqref{eq:student}）を学習する．
事前確率の調整と再均衡化は，出発点の出力を補正する．
先生を必要とする比較手法には，融合の出発点（単一モダリティの生徒に対する温度つきの蒸留とハードな擬似ラベル），$K$本のラベルで学び直した融合の生徒（一般化蒸留），単一モダリティの出発点の積（融合の生徒に対する蒸留）のいずれかを先生として与える．
2段階の参照はどの比較手法にも与えないので，差には先生の選び方の違いも含まれる．
融合の生徒に対しては一般化蒸留とハードな擬似ラベルは定義されず，共訓練は2つの変種で加わる．
各比較手法の\emph{調整版}（tuned）は，提案法と同じく，そのハイパーパラメータ（しきい値，温度，ラベルのない損失の重み）を同じ仮のターゲットで$K$ごとに選ぶ．
\emph{既定値版}（fixed）は元の論文の値を使う．
提案法（ours）は\ref{sec:part2-crossmodal}節の基本の手法を，提案法の拡張版（ours-ext）は拡張した候補群を指す．
AVTの構成では音声，顔，テキスト，融合の生徒を\NUM{cmp-targets-avt-n}のターゲットで評価し，AV，VT，ATの構成では融合の生徒をそれぞれ\NUM{cmp-targets-av-n}，\NUM{cmp-targets-vt-n}，\NUM{cmp-targets-at-n}のターゲットで評価し，2つめのAVTの構成では，すべてのエンコーダを第2のエンコーダ（HuBERT，CLIP，DeBERTa）に置き換える．
結果はターゲットごとに\NUM{cmp-reps-n}回の反復の平均であり，仮のターゲットでは\NUM{cmp-pseudo-reps-n}回の反復を用いる．

\subsection{結果}
\label{sec:comparison-results}

提案法の拡張版は，$K=0$と$K=5$のすべての生徒で，調整版の比較手法のうち最も強いものを\NUM{cmp-oursx-vs-best-range-lo}--\NUM{cmp-oursx-vs-best-range-hi}ポイント上回る（表\ref{tab:compare}．ターゲット間の範囲は補足資料の図S6に示す）．
どこでも最も強い比較手法は1つも無い．
先生の候補に$K$本のラベルで学び直した融合の生徒を含まない基本の手法は，生徒と$K \in \{0, 5\}$の\NUM{cmp-ours-vs-best-cells-n}の組み合わせのうち\NUM{cmp-ours-vs-best-pos-n}で最も強い比較手法を上回る．
$K=5$では，AVTの音声とテキストの生徒で一般化蒸留を下回り，ATの融合の生徒で共訓練と並ぶ．
ラベルの本数にわたって見ると（図S7），提案法の拡張版は音声の生徒ではすべての$K$で優位を保つ（\NUM{cmp-oursx-vs-best-avt-a-lo}--\NUM{cmp-oursx-vs-best-avt-a-hi}ポイント）．
一般化蒸留が提案法の拡張版を上回るのは$K \ge 10$のテキストの生徒だけで，その差は\NUM{cmp-gd-vs-oursx-avt-t-kge10-lo}--\NUM{cmp-gd-vs-oursx-avt-t-kge10-hi}ポイントである．
5つの値を選んだ出発点に対して，提案法の拡張版は，AVTの音声の生徒ではすべての$K$で\NUM{cmp-oursx-vs-e39-avt-a-lo}--\NUM{cmp-oursx-vs-e39-avt-a-hi}ポイント，テキストでは$K \le 5$で\NUM{cmp-oursx-vs-e39-avt-t-kle5-lo}--\NUM{cmp-oursx-vs-e39-avt-t-kle5-hi}ポイント，融合の生徒では$K \ge 1$で\NUM{cmp-oursx-vs-e39-avt-f-kge1-lo}--\NUM{cmp-oursx-vs-e39-avt-f-kge1-hi}ポイントの利得を得て，基本の手法に見られた$K \ge 10$での損失は生じない．
各差は，ターゲットと反復の組にわたる両側のWilcoxonの符号順位検定で検定し，すべての生徒とラベルの本数にわたってHolmの方法で多重性を補正する．
最も強い比較手法に対する差は，$K=0$では\NUM{st-a-k0-cells-n}の生徒のうち\NUM{st-a-k0-pos-n}で，$K=5$では\NUM{st-a-k5-cells-n}のうち\NUM{st-a-k5-pos-n}で正であり，効果量（effect size）$d_z$（差の平均をその標準偏差で割った値）は\NUM{st-a-k0-dz-lo}--\NUM{st-a-k0-dz-hi}と\NUM{st-a-k5-dz-lo}--\NUM{st-a-k5-dz-hi}である．
$5\%$水準で有意なのはそれぞれ\NUM{st-a-k0-sig-n}と\NUM{st-a-k5-sig-n}の生徒である．
1つのターゲットのいくつかの反復はテストの分割を共有するので，分割を共有する反復を平均したうえでも検定する．
このとき有意な生徒はそれぞれ\NUM{st-a-k0-sigdedup-n}と\NUM{st-a-k5-sigdedup-n}になる（$K=0$ではAVTの音声の生徒とVTとATの融合の生徒，$K=5$ではAVTの顔とテキストの生徒とAVとATの融合の生徒を除くすべて）．
構成ごとのターゲットは4から6つなので，これらの検定は本論文のコーパスの中のばらつきを記述するものであり，他のコーパスへの一般化を示すものではない．

\begin{table*}[t]
\caption{1つの手順の下での比較．5つの値を選んだ出発点のUARと，調整版の最も強い比較手法，提案法，提案法の拡張版の出発点からの差（ポイント）を，ターゲットと10回の反復にわたって平均して示す．「なし」は，調整版の比較手法がすべて，ラベルのないデータを使わない設定を選んだことを表す．比較手法の略号は次のとおりである．GD：一般化蒸留，KD：温度つきの蒸留，HPL：ハードな擬似ラベル，RB：出力の再均衡化，ST：自己訓練，CoT：共訓練，CoR：共正則化．}
\label{tab:compare}
\centering
\footnotesize
\setlength{\tabcolsep}{4pt}
\begin{tabular}{@{}llrrlrrrrrr@{}}
\toprule
 & & \multicolumn{2}{c}{出発点} & 最強の & \multicolumn{2}{c}{最強} & \multicolumn{2}{c}{提案法} & \multicolumn{2}{c}{提案法の拡張版} \\
 & & \multicolumn{2}{c}{（UAR）} & 比較手法 & \multicolumn{2}{c}{$-$ 出発点} & \multicolumn{2}{c}{$-$ 出発点} & \multicolumn{2}{c}{$-$ 出発点} \\
\cmidrule(lr){3-4}\cmidrule(lr){6-7}\cmidrule(lr){8-9}\cmidrule(l){10-11}
構成 & 生徒 & $K{=}0$ & $K{=}5$ & （$K{=}0$/$5$） & $K{=}0$ & $K{=}5$ & $K{=}0$ & $K{=}5$ & $K{=}0$ & $K{=}5$ \\
\midrule
AVT & A & \NUM{tab-avt-a-k0-e39} & \NUM{tab-avt-a-k5-e39} & GD/GD & \NUM{cmp-best-vs-e39-avt-a-k0} & \NUM{cmp-best-vs-e39-avt-a-k5} & \NUM{cmp-ours-vs-e39-avt-a-k0} & \NUM{cmp-ours-vs-e39-avt-a-k5} & \NUM{cmp-oursx-vs-e39-avt-a-k0} & \NUM{cmp-oursx-vs-e39-avt-a-k5} \\
 & V & \NUM{tab-avt-v-k0-e39} & \NUM{tab-avt-v-k5-e39} & RB/なし & \NUM{cmp-best-vs-e39-avt-v-k0} & \NUM{cmp-best-vs-e39-avt-v-k5} & \NUM{cmp-ours-vs-e39-avt-v-k0} & \NUM{cmp-ours-vs-e39-avt-v-k5} & \NUM{cmp-oursx-vs-e39-avt-v-k0} & \NUM{cmp-oursx-vs-e39-avt-v-k5} \\
 & T & \NUM{tab-avt-t-k0-e39} & \NUM{tab-avt-t-k5-e39} & HPL/GD & \NUM{cmp-best-vs-e39-avt-t-k0} & \NUM{cmp-best-vs-e39-avt-t-k5} & \NUM{cmp-ours-vs-e39-avt-t-k0} & \NUM{cmp-ours-vs-e39-avt-t-k5} & \NUM{cmp-oursx-vs-e39-avt-t-k0} & \NUM{cmp-oursx-vs-e39-avt-t-k5} \\
 & 融合 & \NUM{tab-avt-f-k0-e39} & \NUM{tab-avt-f-k5-e39} & RB/ST & \NUM{cmp-best-vs-e39-avt-f-k0} & \NUM{cmp-best-vs-e39-avt-f-k5} & \NUM{cmp-ours-vs-e39-avt-f-k0} & \NUM{cmp-ours-vs-e39-avt-f-k5} & \NUM{cmp-oursx-vs-e39-avt-f-k0} & \NUM{cmp-oursx-vs-e39-avt-f-k5} \\
AV & 融合 & \NUM{tab-av-f-k0-e39} & \NUM{tab-av-f-k5-e39} & RB/ST & \NUM{cmp-best-vs-e39-av-f-k0} & \NUM{cmp-best-vs-e39-av-f-k5} & \NUM{cmp-ours-vs-e39-av-f-k0} & \NUM{cmp-ours-vs-e39-av-f-k5} & \NUM{cmp-oursx-vs-e39-av-f-k0} & \NUM{cmp-oursx-vs-e39-av-f-k5} \\
VT & 融合 & \NUM{tab-vt-f-k0-e39} & \NUM{tab-vt-f-k5-e39} & RB/ST & \NUM{cmp-best-vs-e39-vt-f-k0} & \NUM{cmp-best-vs-e39-vt-f-k5} & \NUM{cmp-ours-vs-e39-vt-f-k0} & \NUM{cmp-ours-vs-e39-vt-f-k5} & \NUM{cmp-oursx-vs-e39-vt-f-k0} & \NUM{cmp-oursx-vs-e39-vt-f-k5} \\
AT & 融合 & \NUM{tab-at-f-k0-e39} & \NUM{tab-at-f-k5-e39} & RB/CoT & \NUM{cmp-best-vs-e39-at-f-k0} & \NUM{cmp-best-vs-e39-at-f-k5} & \NUM{cmp-ours-vs-e39-at-f-k0} & \NUM{cmp-ours-vs-e39-at-f-k5} & \NUM{cmp-oursx-vs-e39-at-f-k0} & \NUM{cmp-oursx-vs-e39-at-f-k5} \\
\midrule
AVT， & A & \NUM{tab-w2-a-k0-e39} & \NUM{tab-w2-a-k5-e39} & CoT/GD & \NUM{cmp-best-vs-e39-w2-a-k0} & \NUM{cmp-best-vs-e39-w2-a-k5} & \NUM{cmp-ours-vs-e39-w2-a-k0} & \NUM{cmp-ours-vs-e39-w2-a-k5} & \NUM{cmp-oursx-vs-e39-w2-a-k0} & \NUM{cmp-oursx-vs-e39-w2-a-k5} \\
第2の & V & \NUM{tab-w2-v-k0-e39} & \NUM{tab-w2-v-k5-e39} & CoT/KD & \NUM{cmp-best-vs-e39-w2-v-k0} & \NUM{cmp-best-vs-e39-w2-v-k5} & \NUM{cmp-ours-vs-e39-w2-v-k0} & \NUM{cmp-ours-vs-e39-w2-v-k5} & \NUM{cmp-oursx-vs-e39-w2-v-k0} & \NUM{cmp-oursx-vs-e39-w2-v-k5} \\
エンコーダ & T & \NUM{tab-w2-t-k0-e39} & \NUM{tab-w2-t-k5-e39} & RB/GD & \NUM{cmp-best-vs-e39-w2-t-k0} & \NUM{cmp-best-vs-e39-w2-t-k5} & \NUM{cmp-ours-vs-e39-w2-t-k0} & \NUM{cmp-ours-vs-e39-w2-t-k5} & \NUM{cmp-oursx-vs-e39-w2-t-k0} & \NUM{cmp-oursx-vs-e39-w2-t-k5} \\
 & 融合 & \NUM{tab-w2-f-k0-e39} & \NUM{tab-w2-f-k5-e39} & RB/CoT & \NUM{cmp-best-vs-e39-w2-f-k0} & \NUM{cmp-best-vs-e39-w2-f-k5} & \NUM{cmp-ours-vs-e39-w2-f-k0} & \NUM{cmp-ours-vs-e39-w2-f-k5} & \NUM{cmp-oursx-vs-e39-w2-f-k0} & \NUM{cmp-oursx-vs-e39-w2-f-k5} \\
\bottomrule
\end{tabular}
\end{table*}

\subsection{既定値のハイパーパラメータとエントロピー最小化}
\label{sec:comparison-fixed}

元の論文のハイパーパラメータを使うと，エントロピー最小化，共正則化，事前確率の調整は，$K=0$でAVTの単一モダリティの生徒のUARを\NUM{cmp-fixed-collapse-avt-mod-k0-lo}から\NUM{cmp-fixed-collapse-avt-mod-k0-hi}ポイント変え，共正則化と，音声の生徒に対するエントロピー最小化は，UARを偶然の水準まで下げる．
融合の生徒では，既定値版のエントロピー最小化は，4つの構成にわたる平均で出発点のUARを\NUM{cmp-entmin-fixed-vs-e39-f-k0-lo}から\NUM{cmp-entmin-fixed-vs-e39-f-k0-hi}ポイント変え，その効果はターゲットによって異なる（図S8）．
IEMOCAPの即興のセッションでは，UARを$+$\NUM{cmp-entmin-fixed-vs-e39-avt-f-k0-impro}（AVT）と$+$\NUM{cmp-entmin-fixed-vs-e39-at-f-k0-impro}ポイント（AT）変えるが，同じターゲットでも\NUM{cmp-entmin-fixed-vs-e39-vt-f-k0-impro}（VT）と\NUM{cmp-entmin-fixed-vs-e39-av-f-k0-impro}ポイント（AV）変え，AVTのMELDでは\NUM{cmp-entmin-fixed-vs-e39-avt-f-k0-meld}ポイント変える．
仮のターゲットによる選択はこれらの失敗を検出し，すべての$K$でAVTの融合の生徒のエントロピー最小化を止める．

\subsection{第2のエンコーダ，微調整，選択}
\label{sec:comparison-robust}

第2のエンコーダでは，基本の手法は$K=0$と$K=5$ですべての生徒を出発点より上に上げ（表\ref{tab:compare}），提案法の拡張版は，調整版の最も強い比較手法を，$K=0$と$K=5$で\NUM{cmp-oursx-vs-best-w2-k0k5-lo}--\NUM{cmp-oursx-vs-best-w2-k0k5-hi}ポイント，$K=25$で\NUM{cmp-oursx-vs-best-w2-k25-lo}--\NUM{cmp-oursx-vs-best-w2-k25-hi}ポイント上回る．

凍結した設定が適応の能力を過小に示していないかを確かめるために，4つのターゲットの音声の生徒について，3回の反復で，WavLMを低ランクのアダプタで微調整した（補足資料の図S9とS7節）．
感情ごとに$K=5$本のラベルで微調整すると，凍結した5つの値を選んだ出発点を\NUM{cmp-ft-konly-below-frozen-e39-k5}ポイント下回る（ターゲット間で\NUM{cmp-ft-konly-below-frozen-e39-k5-tgt-lo}--\NUM{cmp-ft-konly-below-frozen-e39-k5-tgt-hi}）．
AVTの音声の生徒で選ばれた形に固定した提案法の拡張版の弱いラベルは，微調整したモデルでも$K=0$で\NUM{cmp-ft-weak-vs-start-k0}ポイント，$K=5$で\NUM{cmp-ft-weak-vs-anchor-k5}ポイント役立つが，提案法の微調整版は，凍結した特徴の上の同じ設定を\NUM{cmp-ft-ours-below-frozen-ours-k0}ポイントと\NUM{cmp-ft-ours-below-frozen-ours-k5}ポイント下回り，中心化はその差の\NUM{cmp-ftc-gap-closed-start-k0}--\NUM{cmp-ftc-gap-closed-konly-k5}\%しか縮めない．

\emph{選択による損失．}
\label{sec:comparison-selection}
仮のターゲットによる選択は，提案法の候補群のうち最良の候補を常に見つけるわけではない．
$K=0$では，テスト話者での最良の候補は選んだ候補を，AVTのテキストの生徒で\NUM{cmp-oursx-oracle-vs-tuned-avt-t-k0}ポイント，第2のエンコーダのテキストの生徒で\NUM{cmp-oursx-oracle-vs-tuned-w2-t-k0}ポイント，その他の生徒で\NUM{cmp-oursx-oracle-vs-tuned-at-f-k0}--\NUM{cmp-oursx-oracle-vs-tuned-w2-v-k0}ポイント上回る．
差が1ポイントに達するのは，\NUM{cmp-oursx-oracle-gap-students-n}の生徒のうち\NUM{cmp-oursx-oracle-gap-ge1-n}である．
この差はよりよい選択の規則が加えうるものの上限を与え，比較手法も同じ手順で選ぶので，提案法に有利には働かない．

\section{考察}
\label{sec:discussion}

\emph{実用上の指針．}
結果から，感情ごとに少数のラベルしかない新しいコーパスに対する4つの推奨が得られる．
第1に，ソースとターゲットの特徴を別々に中心化する．
これはラベルを使わず，ソースのみの分類器がもつ判断基準による損失の大部分を取り除く（\ref{sec:part1-fix}節）．
第2に，すべてのソースから出発し，重みには強いアンカーをかけたまま，緩いアンカーをかけた切片に$K$本のラベルを使う．
ラベルは切片をよく決め，重みをうまく決めないからである（\ref{sec:part2-start}節）．
第3に，どの発話にラベルを付けるかをラベルを使わずに選んでも，よい出発点の向きが直ることは期待しない（\ref{sec:part2-limits}節）．
第4に，ターゲットの発話とともに他のモダリティが記録されているならば，それがラベル付けのときだけであっても，それらのモダリティでソースから学習した先生にプール全体のラベルを付けさせる．
配備する分類器は単一モダリティのままでよい（\ref{sec:part2-crossmodal}節と\ref{sec:comparison}節）．

\emph{テキストのふるまいが異なる理由．}
テキストの特徴は判断基準による損失が小さく，向きによる損失が大きい（表\ref{tab:part1}）．
このため，他の構成で判断基準を直す少数のラベルには，テキストでは直すべき判断基準のずれがほとんど無く，テキストだけの構成では利得のうち判断基準の割合は最大で\NUM{p1-e29-critshare-k1-t}である．
2つのコストによる同じ説明がこの場合にも当てはまる．
利得のうち判断基準から生じる割合は，損失のうち判断基準が引き起こす割合に従う．
提案法では，テキストの生徒も音声の生徒と同じくらい融合の先生から利得を得る．
融合の先生が，テキストのラベルに欠けている向きの情報を補うからである．

\emph{エントロピー最小化とクラスタの仮定．}
エントロピー最小化は判定の境界をラベルのないデータの密な領域から遠ざけ，クラスが分かれたクラスタをなすときに役立つ．
本論文のデータでは，エントロピー最小化は，AVTとATの構成のIEMOCAPの即興のセッションでは役立ち，MELDとCMU-MOSEI，およびVTとAVの構成の同じIEMOCAPのターゲットでは大きく失敗する（\ref{sec:comparison-fixed}節）．
別のモダリティからの弱いラベルは，モダリティ間の相補性という別の性質に依拠しており，その利得のばらつきはずっと小さい．
$K=0$と$K=5$のターゲットとセルの\NUM{st-b-tgt-n}の組み合わせのうち，提案法の拡張版が出発点を下回るのは\NUM{st-b-tgtneg-n}で，その差は最大で\NUM{st-b-tgtneg-worst}ポイントである．

\emph{ソースとは何か．}
第I部はソースの価値（value of a source）を向きのラベルの本数で与え，第II部は先生の価値が出発点に対する先生のUARの優位に応じて増減することを示す．
どちらの量も配備の前に推定できる．
$q$と$e$はラベル付きのソースと仮のターゲットから，$J$は仮のターゲットから推定する．
これらの量は，ターゲットのラベルを増やすこと，ソースを加えること，ラベル付けのときにモダリティを加えることのどれがよりよい投資かを示す．

\section{限界と倫理}
\label{sec:limitations-ethics}

\emph{分析の範囲．}
すべての分析は，凍結したエンコーダと線形ヘッドを用いる．
微調整の確認は，低ランクのアダプタで適応した1つの音声エンコーダを4つのターゲットと3回の反復で調べたものなので，主張を限定できるのはその設定についてだけである．
2つのガウス分布のモデルは，各感情の対の中でスコア差が共通の分散をもつガウス分布に従うと仮定し，式\eqref{eq:lambda}，式\eqref{eq:value}，命題\ref{prop:design}は1次または2次の近似に依拠する．
本論文は厳密な結果を主張せず，それらの予言を実験で検証する．
判断基準が先に直るという説明が成り立つのは，単一モダリティの音声と顔の特徴で$K \le 5$の場合である．
融合の構成とテキストの構成では，$K=5$で判断基準の部分は利得のおよそ半分かそれ以下である．

\emph{データ．}
6つのコーパスから作った7つの群の4つの感情を調べ，すべて英語である．
コーパスごとにラベルの定義が異なり，特にCMU-MOSEIの中立のクラスは評定された感情が無いことで定義される．
CMU-MOSEIではほとんどの動画のチャンネルが分からないので話者の分離は不完全であり，RAVDESSには本論文の処理の流れで顔の特徴が無い．
テキストを含む構成には4つの群しかなく，手法の比較ではテキストを含むどの構成でも4つのターゲットを使うので，ターゲット間の違いは少数のターゲットから推定している．
第2のエンコーダは，提案法については3つのモダリティすべてを含む構成でだけ検証した．

\emph{統計と選択．}
本論文の区間はターゲットのラベルとテスト話者の引き方による変動を表し，話者やコーパスの再標本化による変動は表さない．
また，IEMOCAPのターゲットでは反復は独立でない．
仮のターゲットでハイパーパラメータを選ぶと，$K=0$ではテスト話者での最良の値に対して最大で\NUM{cmp-oursx-oracle-vs-tuned-avt-t-k0}ポイントを失う（\ref{sec:comparison-selection}節）．
列の正規化は，本論文のターゲットのプールで感情の数が偏っているにもかかわらず（表\ref{tab:data}．たとえばMELD），感情ごとの合計をそろえる．
列の正規化による利得がクラスの比率にどう依存するかは分析しておらず，偏りがさらに強いと列の正規化は害になりうる\cite{veilleux2021realistic}．

\emph{倫理．}
公開されたコーパスだけをその配布条件の下で使い，新たなデータは記録していない．
顔は，検出した顔の切り出しの埋め込みとしてだけ処理の流れに入り，個人を識別できるデータは公開しない．
感情認識は，監視や，本人の同意のない人についての決定に悪用されうる．
本論文の手法は演技とメディアの録画からカテゴリのラベルを推定するものであり，個人についての決定のためには検証されていない．

\section{結論}
\label{sec:conclusion}

感情ごとに少数のターゲットのラベルがあり，凍結した特徴の上で線形ヘッドを使うとき，ラベルはまず判定の判断基準を直す．
単一モダリティの音声と顔で$K \le 5$のとき，ラベルが合併の分類器に加える改善の\NUM{p1-e29-critshare-kle5-single-lo}--\NUM{p1-e29-critshare-kle5-single-hi}は判断基準から生じ，一方で音声の分類器の向きは$q/K$の速さでしか改善しない（$q$は\NUM{p1-qdir-audio-lo}から\NUM{p1-qdir-audio-hi}）．
音声では，ソースは向きについて感情ごとに\NUM{p1-e32-n1src-audio-lo}--\NUM{p1-e32-n1src-audio-hi}本のラベルに相当する情報を補い，少数のラベルも，ラベルを付ける発話をラベルを使わずに選ぶことも，よい出発点の向きをあまり動かさない．
ソースだけで学習し，生徒がもたないモダリティを見る融合の先生がラベルのないプール全体に付ける弱いラベルは，その向きを動かす．
$K \le 5$では，この弱いラベルはすべての単一モダリティの生徒を5つの値を選んだ出発点より上に上げ，ターゲットのプールで学び直した先生を加えた拡張版は，$K=0$と$K=5$のすべての生徒について，第2のエンコーダでも，同じ手順で調整した最も強い比較手法を上回る．
$q$，$e$，出発点に対する先生のUARの優位$J$はターゲットのラベルを使って測ったので，ラベルを増やすこと，別のソースを加えること，別のモダリティを加えることのどれを選ぶかをラベル付けの前に決めるために，これらを仮のターゲットで推定することが次の課題である．


\bibliographystyle{../IEEEtran}
\bibliography{../references}

\end{document}
